% Diagrama: Arquitectura hexagonal de tres capas de VFTModel
% Usado en: 4-Capas-Codigo/4-2-VFTModel.tex  (fig:vft-arquitectura)
\begin{tikzpicture}[
    capa/.style = {
        rectangle, draw=unamAzul, rounded corners=5pt,
        text width=11cm, align=left,
        minimum height=1.6cm, inner sep=10pt
    },
    flecha/.style = {Stealth-Stealth, thick, color=unamAzul!55}
]
    \node[capa, fill=unamAzul!5] (api) {
        \textbf{Capa de \api{}} \hfill \texttt{src/api/}\\
        {\small Peticiones \textsc{http} · Validación Pydantic · Respuestas \textsc{json}}
    };
    \node[capa, fill=unamAzul!12, below=0.45cm of api] (core) {
        \textbf{Capa de Dominio} \hfill \texttt{src/core/}\\
        {\small Algoritmos topológicos · Modelos matemáticos · Constructor del grafo}
    };
    \node[capa, fill=unamAzul!22, below=0.45cm of core] (infra) {
        \textbf{Capa de Infraestructura} \hfill \texttt{src/infrastructure/}\\
        {\small Clientes \textsc{http} · Archivos \geojson{} locales · Colas asíncronas}
    };
    \draw[flecha] (api)  -- (core);
    \draw[flecha] (core) -- (infra);
\end{tikzpicture}
